\section{Connectomics and Volumetric Imaging}

Connectomics is the study of neural wiring diagrams at synaptic resolution.
A connectome maps which neurons connect to which, with what synaptic strength,
and through which anatomical pathways. To obtain this information, researchers
section brain tissue into ultrathin slices (30--50~nm thick for EM; thicker
sections are used for XNH), image each slice, and reconstruct the
three-dimensional morphology of every neuron computationally.

Two imaging modalities are relevant to this work:

\paragraph{Electron Microscopy (EM).}
Serial-section transmission EM and focused ion beam scanning EM (FIB-SEM)
achieve 4--8~nm lateral resolution, sufficient to resolve individual synaptic
vesicles. The CREMI dataset~\cite{cremi2016} used in this work is a
$1250 \times 1250 \times 125$ voxel volume of \emph{Drosophila} larval brain
at $4 \times 4 \times 40$~nm resolution. The density of membranes and
organelles in EM images makes segmentation challenging: membranes appear as
bright thin sheets that segmenters must trace accurately.

\paragraph{X-ray Nano-Holotomography (XNH).}
XNH uses synchrotron X-rays to image brain tissue in three dimensions without
physical sectioning. It achieves lower resolution than EM (typically 25--33~nm
isotropic) but can image much larger volumes non-destructively. The XPRESS
challenge~\cite{xpress2023} uses a $699 \times 699 \times 699$ voxel volume of
mouse white matter at 33~nm isotropic resolution. Myelinated axons appear as
bright tubular structures surrounded by darker myelin sheaths, giving them
a characteristic ring-like cross-section that aids geometric reasoning.

\section{Automated Neural Segmentation}

Modern neural segmenters use deep convolutional networks, typically operating
on affinity predictions (per-voxel estimates of whether neighboring voxels
belong to the same object). Affinity maps are then processed by a watershed
or flood-filling algorithm to produce a label volume where each integer value
identifies a unique putative neuron fragment.

Flood-filling networks~\cite{januszewski2018high} are among the most accurate
segmenters: a convolutional network iteratively ``fills'' a segment from a seed,
deciding at each step whether to extend the current object or stop. Despite
high accuracy, all automated approaches produce residual errors.

\section{Segmentation Errors: Splits and Merges}

Segmentation errors fall into two categories:

\begin{description}
  \item[Merge errors] Two distinct neurons are assigned the same label. A merge
        introduces a false synapse (two neurons that appear connected when they
        are not) and is considered the more serious error type.
  \item[Split errors] A single neuron is broken into multiple fragments with
        distinct labels. Splits are more common than merges in long, thin structures
        like axons, where slight imaging artifacts or membrane contrast variation
        can interrupt a segmenter's flood-fill.
\end{description}

The XPRESS challenge reports that the baseline segmenter produces 1{,}499
ground-truth merge pairs across the training volume, with each pair representing
one split error that the pipeline must detect and propose to merge. The challenge
is to accept as many true pairs as possible (high recall) while rejecting
cross-axon false positives (high precision).

\section{Skeletonization}

Skeleton extraction converts a volumetric label into a compact graph of
centerline nodes and edges, capturing the topology of the neuron's morphology.
The TEASAR algorithm (Tree-structure Extraction Algorithm for Accurate and Robust
Skeletons)~\cite{sato2000teasar} traces a tree-structured skeleton by finding
shortest paths to a set of ``target'' voxels identified via distance-transform
analysis. TEASAR produces compact, well-centered skeletons with configurable
branch sensitivity. We use the \texttt{kimimaro}~\cite{kimimaro} Python
implementation, which supports GPU acceleration and multi-label batch processing.

Skeleton endpoints (degree-1 nodes in the skeleton graph) are natural
connection points: a true split typically produces two fragments whose free
endpoints point toward each other across the gap. However, for long axons that
are split in the interior (far from either end), the relevant skeleton nodes are
\emph{interior} nodes, not endpoints. This is a key observation that motivates the
skeleton-node graph construction described in Chapter~\ref{ch:methodology}.

\section{The XPRESS Challenge}

The XPRESS (X-ray nano-holotomography Processing for Evaluation of Segmentation
Software) challenge~\cite{xpress2023} provides:

\begin{itemize}
  \item A baseline segmentation volume (\texttt{baseline\_seg\_full.h5}): the
        output of an automated segmenter applied to mouse white matter XNH data.
        The volume is $699^3$ voxels at 33~nm isotropic resolution.
  \item A ground-truth skeleton dataset (\texttt{XPRESS\_training\_skels.npz}):
        manually traced skeleton graphs for the true axon trajectories, used as
        a proofreading oracle.
  \item A held-out validation volume (\texttt{baseline\_seg\_validation.h5}) with
        203 ground-truth merge pairs, used to evaluate generalization.
\end{itemize}

Ground-truth evaluation proceeds as follows: for each true axon trajectory, the
two baseline-segmentation fragments whose label IDs contain the most ground-truth
skeleton nodes at either end of the trajectory define an oracle merge pair. A
pipeline \emph{covers} a pair if it generates at least one candidate connecting
those two fragments. Of the covered pairs, those that are also accepted constitute
true positives.

\section{Prior Approaches to Automated Proofreading}

Early proofreading systems were purely manual, relying on human experts working
in tools such as CATMAID or Neuroglancer to trace axons and correct errors. The
scale of modern datasets has made fully manual proofreading infeasible, motivating
a shift toward human-in-the-loop systems~\cite{dorkenwald2022cave} where
algorithms propose corrections and humans verify only the uncertain cases.

Automated correction proposal systems generally follow one of two paradigms:
\emph{voxel-level} approaches, which re-run or fine-tune segmentation on regions
flagged as likely errors, and \emph{graph-level} approaches, which reason about
fragment-to-fragment merge plausibility using geometric or learned features.
This work falls in the latter category, operating entirely in the post-segmentation
graph domain without re-examining raw image data.
